\section*{Chapitre \ref{ch:qm:numerical-optimisation-in-practice} --- Optimisation numérique en pratique}

\begin{solution}{exo:qm:numerical-optimisation-in-practice:1}
Avec la borne $(1 - 1/\kappa)^k$, $k \approx \kappa\ln10^6 = 13\,816$ itérations~; avec accélération, environ $\sqrt\kappa\ln10^6 = 437$.
\end{solution}

\begin{solution}{exo:qm:numerical-optimisation-in-practice:2}
Seuillage doux de 1~: $(2, 0, 0)$.
\end{solution}

\begin{solution}{exo:qm:numerical-optimisation-in-practice:3}
Échanger $(a, \tau_1, \tau_2)$ contre $(1 - a, \tau_2, \tau_1)$ donne la même fonction du retard~: le modèle est invariant par renommage de ses composantes, de sorte que chaque minimum possède une image miroir.
Ordonner les échelles de temps dans le paramétrage l'élimine.
\end{solution}

\begin{solution}{exo:qm:numerical-optimisation-in-practice:4}
Pour $f(x) = \frac12x^\top Hx + b^\top x$ avec $H$ définie positive, $x - H^{-1}(Hx + b) = -H^{-1}b$, l'unique minimiseur, quel que soit $x$.
\end{solution}

\begin{solution}{exo:qm:numerical-optimisation-in-practice:5}
Avec $s = td$, $y = \nabla f(x + td) - \nabla f(x)$~: $y^\top s = t(\nabla f(x + td)^\top d - \nabla f(x)^\top d) \ge t(c_2 - 1)\nabla f(x)^\top d > 0$, puisque $\nabla f(x)^\top d < 0$ et $c_2 < 1$.
\end{solution}

\begin{solution}{exo:qm:numerical-optimisation-in-practice:6}
$1/\sqrt k$~: la somme diverge, mais la somme des carrés diverge aussi ($\sum1/k$), donc la forme classique n'est pas satisfaite (bien que la moyenne des itérés la rétablisse). $1/k$~: la somme diverge et la somme des carrés
converge, donc oui.
\end{solution}

\begin{solution}{exo:qm:numerical-optimisation-in-practice:7}
\texttt{daily(w)} pour $w = 0$, 0.001, 0.01, 0.1~: variation médiane 4.48, 3.93, 1.87 et 0.42 jours~; RMSE moyenne ($\times 10^{-3}$) 9.450, 9.454, 9.520 et 9.643. La stabilité coûte peu jusqu'à 0.01, puis commence à
coûter en qualité d'ajustement.
\end{solution}

\begin{solution}{exo:qm:numerical-optimisation-in-practice:8}
La convergence vers un point stationnaire ne dit rien de quel minimum a été trouvé (images miroirs, solutions au bord), ni de la platitude de l'objectif en ce point~: les paramètres d'une vallée plate se déplacent
avec le bruit, alors que le gradient est nul au point de chaque jour. Il faut vérifier le \omterm{def:qm:numerical-optimisation-in-practice:calibration}{multistart}, le conditionnement de $J^\top J$ et la stabilité d'un jour à l'autre.
\end{solution}

\begin{solution}{pb:qm:numerical-optimisation-in-practice:1}
\textbf{1.} $a = 0.59$, $\tau_1 = 2.9$ et $\tau_2 = 55$ jours, RMSE 0.0096.
\textbf{2.} 28 au meilleur ajustement, 35 à son image miroir, 26 à la frontière de 1\,000 jours avec au moins le double de l'erreur, 11 ailleurs.
\textbf{3.} $(0.41, 55, 2.9)$~: la même courbe avec les étiquettes des composantes échangées.
\textbf{4.} Conditionnement 10.5~; corrélation 0.96 entre le poids rapide et l'échelle de temps lente.
\textbf{5.} D'environ 45 à 80 jours (erreur de profil $9.9 \times 10^{-3}$ à 62 jours, 12.3 à 42, 12.9 à 89, contre un bruit de 10).
\textbf{6.} Médiane de 4.5 jours, maximum 30.7.
\textbf{7.} Jour 189~: l'échelle de temps lente est passée de 56.8 à 87.6 jours et l'erreur d'ajustement de 0.00926 à 0.00921.
\textbf{8.} Médiane de 1.9 jours, maximum 9.5.
\textbf{9.} La RMSE moyenne passe de 0.00945 à 0.00952, soit 0.7~\%.
\textbf{10.} Son écart-type passe de 5.5 à 2.7 jours.
\textbf{11.} La \omterm{def:qm:numerical-optimisation-in-practice:gd}{descente de gradient} atteint 0.067 après 100 itérations et 0.0090 après 200~; Nesterov $5 \times 10^{-6}$ après 100.
\textbf{12.} Un pas de Newton (deux itérations en comptant la vérification)~; L-BFGS en fait sept.
\textbf{13.} ISTA est à $7.6 \times 10^{-4}$ de l'optimum après 50 itérations~; FISTA environ $10^{-7}$ après 50 et inférieur à $10^{-10}$ après 65.
\textbf{14.} 91 itérations d'ADMM contre 74 projections, la même matrice à $1.4 \times 10^{-9}$ près.
\textbf{15.} Convergence~: à 0.009 de la solution des moindres carrés après 20\,000 pas, contre 0.113 avec un pas constant.
\textbf{16.} L'échelle de temps lente, avec son intervalle le long de la vallée (environ 45 à 80 jours) et sa variation d'un jour à l'autre.
\textbf{17.} Par \omterm{def:qm:linear-models-under-stress:cv}{validation croisée}~: le poids qui minimise l'erreur d'ajustement de demain des paramètres d'aujourd'hui, ou un poids qui plafonne la variation d'un jour à l'autre à une tolérance acceptée par les utilisateurs.
\textbf{18.} Paramétrer $\tau_2 = \tau_1 + e^u$ (et $a \in (0, 1)$), de sorte qu'un seul étiquetage soit représentable.
\textbf{19.} Résultat nommé~: \emph{le lundi où l'ajustement a bougé}~: recalage quotidien à partir d'un point de départ fixe, l'échelle de temps lente saute d'une médiane de 4.5 jours (maximum 30.7)~; avec une pénalité vers la veille (poids
0.01), le saut médian tombe à 1.9 jours (maximum 9.5), au prix de 0.7~\% d'erreur d'ajustement.
\textbf{20.} Lequel parmi plusieurs jeux de paramètres aussi bons il a trouvé, et avec quelle précision les données les déterminent.
\end{solution}

\begin{solution}{iq:qm:numerical-optimisation-in-practice:1}
Le pas doit être assez petit pour la direction la plus raide ($1/L$), de sorte que le progrès selon la direction la plus plate ne vaut que $\mu/L$ par pas~: la vitesse est fixée par le conditionnement. Remèdes~: remettre à l'échelle les
variables, préconditionner, exploiter la courbure (Newton, quasi-Newton, Gauss--Newton) ou accélérer (Nesterov, \omterm{def:qm:numerical-optimisation-in-practice:momentum}{méthode du moment}).
\iqlookfor{Conditionnement comme vitesse, et au moins deux remèdes.}
\end{solution}

\begin{solution}{iq:qm:numerical-optimisation-in-practice:2}
L'objectif est plat selon certaines directions (un problème d'\omterm{def:qm:numerical-optimisation-in-practice:calibration}{identifiabilité}), ou présente plusieurs minima équivalents, et le bruit déplace la réponse. Diagnostiquer avec le conditionnement de $J^\top J$, un profil de
l'objectif et le \omterm{def:qm:numerical-optimisation-in-practice:calibration}{multistart}~; corriger par une pénalité vers la veille, un reparamétrage, la fixation d'un paramètre mal identifié, ou des données plus informatives.
\iqlookfor{Diagnostic d'\omterm{def:qm:numerical-optimisation-in-practice:calibration}{identifiabilité}, et une stabilisation qui coûte peu en ajustement.}
\end{solution}

\begin{solution}{iq:qm:numerical-optimisation-in-practice:3}
Il résout $(J^\top J + \lambda D)\Delta x = -J^\top r$~: un $\lambda$ petit donne Gauss--Newton (rapide près de la solution), un $\lambda$ grand un court pas de gradient mis à l'échelle (sûr loin d'elle)~; $\lambda$ s'adapte au
succès. Gauss--Newton pur peut faire d'énormes pas lorsque $J^\top J$ est presque singulière, précisément dans les vallées plates.
\iqlookfor{Interpolation entre Gauss--Newton et \omterm{def:qm:numerical-optimisation-in-practice:gd}{descente de gradient}, et la matrice $J^\top J$ singulière.}
\end{solution}

\begin{solution}{iq:qm:numerical-optimisation-in-practice:4}
Chaque gradient stochastique a une variance qui ne s'annule pas à l'optimum, de sorte qu'un pas constant maintient l'itéré en rebond dans une boule de bruit de rayon proportionnel au pas. Avec $\sum t_k = \infty$
(déplacement total suffisant) et $\sum t_k^2 < \infty$ (bruit qui s'éteint), les itérés convergent.
\iqlookfor{Bruit à l'optimum et les deux conditions.}
\end{solution}

\begin{solution}{iq:qm:numerical-optimisation-in-practice:5}
$\operatorname{prox}_{tg}(v)$ minimise $g(x) + \lVert x - v\rVert^2/2t$~; pour la norme $\ell_1$, c'est le seuillage doux. La pénalité non lisse du \omterm{def:qm:linear-models-under-stress:penalties}{lasso} est traitée exactement par l'\omterm{def:qm:numerical-optimisation-in-practice:prox}{opérateur proximal}, tandis que la partie lisse
des moindres carrés utilise son gradient~: la \omterm{def:qm:numerical-optimisation-in-practice:prox}{méthode du gradient proximal} (ISTA, FISTA) et la descente de coordonnées l'exploitent toutes deux.
\iqlookfor{Seuillage doux et décomposition en partie lisse et partie simple.}
\end{solution}

\begin{solution}{iq:qm:numerical-optimisation-in-practice:6}
La mise à jour BFGS ne garde l'approximation de l'inverse de la hessienne définie positive que si $y^\top s > 0$~; la condition de courbure de Wolfe le garantit, alors qu'une simple condition de décroissance ne le fait pas. Sans elle
la direction suivante peut ne pas être une direction de descente.
\iqlookfor{La condition de courbure entraîne la définie-positivité.}
\end{solution}
